\documentclass{article}
\usepackage{neurips_2026}
\usepackage{amsthm}
\usepackage{natbib}
\usepackage{times}
\usepackage{helvet}
\usepackage{courier}
\usepackage{algorithm}
\usepackage{algorithmic}
\usepackage{caption}
\usepackage{subcaption}
\usepackage{enumitem}

\usepackage[utf8]{inputenc}
\usepackage[T1]{fontenc}
\usepackage{hyperref}
\usepackage{url}
\usepackage{booktabs}
\usepackage{amsfonts}
\usepackage{nicefrac}
\usepackage{microtype}
\usepackage{graphicx}
\usepackage{multirow}
\usepackage{amsmath}
\usepackage{amssymb}
\usepackage{mathtools}
\usepackage{xcolor}
\usepackage{colortbl}
\usepackage{cite}

\title{The Truncation Signature: Free Signals and Deterministic Heuristics for Efficient Test-Time Compute Scaling}

\author{
  Anonymous Authors \\
  Anonymous Institution \\
  \texttt{anonymous@example.com}
}

\begin{document}

\maketitle

\begin{abstract}
Output token count is a free correctness signal for test-time compute allocation that requires no model internals (AUC $0.865$ vs.\ $0.561$ for log-probability on 7B models). We show this signal works because model failures at a given budget split into two types with fundamentally different properties: \textbf{Type~A} (truncation: detectable via token count $\approx$ budget, recoverable 0--80\% with more tokens) and \textbf{Type~B} (reasoning: output complete but wrong, undetectable, irrecoverable). Under deterministic decoding, Type~B non-recovery is a mathematical certainty of autoregressive generation---$0/521$ natural-stop errors recover across 3 model families and 3 tasks (a sampling validation at temperature $= 0.7$ confirms the practical asymmetry extends beyond greedy decoding)---and we show this certainty yields a practical \emph{zero-overhead} budget allocation heuristic: extend budget only for truncated outputs.

A Three-Zone heuristic based on ceiling rate (TRUNCATION $>$50\%: extend; INFLECTION 10--50\%: cautious extension; EFFICIENT $<$10\%: stop) requires no verifiers, training, or ground truth. The rule extends to reasoning models: DeepSeek-R1-Distill-7B shows a $2.0\times$ right-shift in the truncation boundary but reaches the same peak accuracy (80.5--81.0\% on MATH/GSM8K), suggesting that budget under-allocation warrants attention at this scale. All experiments use greedy decoding with a sampling validation at temperature $= 0.7$.
\end{abstract}

\section{Introduction}

How much compute should we allocate at inference time? Test-time compute (TTC) scaling has become central to LLM deployment, with approaches ranging from extended chain-of-thought~\cite{wei2022chain} to repeated sampling~\cite{brown2024large,snell2024scaling} to reasoning model distillation~\cite{guo2025deepseek}. A growing literature studies when to stop reasoning---recent work shows that models often ``decide'' answers early and generate post-decision explanation~\cite{decideearly2026}, that reflection after an initial answer provides diminishing returns~\cite{firsttry2025}, and that overthinking can actively hurt accuracy~\cite{morethinkinghurts2026}. Yet these works focus on \emph{when} reasoning becomes wasteful; they do not provide a \emph{decision rule} for budget allocation that is both free (no model internals) and universal (all model types). We ask: is there a single, trivially observable signal that solves the complete budget allocation problem?

We show that there is: \textbf{output token count}. On 7B models, checking whether the output length equals the budget detects wrong answers with AUC $0.865$, compared to $0.561$ for log-probability. This free signal works because model failures split into two types with radically different properties:

\textbf{Layer 1: The Practical Signal---Token Count as Correctness Detector.}
Output token count achieves AUC $0.865$ for detecting wrong answers on Qwen2.5-7B, compared to $0.561$ for log-probability (near-random). The signal is free---requiring only output length, no model internals---and cross-family: Gemma-2-9B and LLaMA-3-8B show AUC $0.699$--$0.731$. Prior work analyzed output length effects~\cite{concisethoughts2024} but without distinguishing truncation from reasoning failures; we show this distinction is what makes the signal work.

\textbf{Layer 2: The Explanation---Two Failure Types.}
The token signal works because failures at a given budget split into two types:
\begin{itemize}
    \item \textbf{Type~A} (Truncation failure): Output was cut off by the budget. Detectable (token count $\approx$ budget, AUC $\approx 1.0$), recoverable (0--80\%, model-dependent).
    \item \textbf{Type~B} (Reasoning failure): Model stopped naturally at wrong answer. Undetectable (AUC $\approx 0.5$), \emph{irrecoverable}.
\end{itemize}
The asymmetry in recoverability is what makes the signal practical: Type~A errors dominate at typical budgets (and are detectable), while Type~B errors never recover regardless (so detecting them provides no actionable information).

\textbf{Layer 3: The Decision Heuristic---Three-Zone Budget Allocation.}
A three-zone heuristic based on ceiling rate---TRUNCATION ($>$50\%: extend), INFLECTION (10--50\%: cautious extension), EFFICIENT ($<$10\%: stop)---requires no verifiers, training, or ground truth. Across our data, this heuristic produces consistent zone assignments across datasets and models.

\textbf{Layer 4: The Foundation---Deterministic Non-Recovery.}
Under deterministic (greedy) decoding, a model that stops naturally at a wrong answer produces the identical output at any higher budget. Type~B non-recovery is therefore a mathematical certainty of autoregressive generation, not an empirical estimate---the novelty lies not in this qualitative observation (which is expected) but in its practical consequence: a zero-overhead heuristic for budget allocation. We verify this across $521$ independent samples (0 recoveries) spanning 3 model families and 3 tasks (math, arithmetic, code). A sampling experiment at temperature $= 0.7$ validates the practical asymmetry extends beyond greedy decoding. We call the combination of deterministic Type~B non-recovery and empirical Type~A recovery the \emph{Ratchet Effect}.

\textbf{Layer 5: The Extension---Reasoning Models.}
We apply this framework to DeepSeek-R1-Distill-7B across 6 budgets (256--8192 tokens) on MATH and GSM8K. The reasoning model requires $2.0\times$ more tokens to escape the truncation regime but reaches the same accuracy ceiling (80.5--81.0\%). While we study only one distilled reasoning model, this provides initial evidence that the framework applies beyond standard instruction-tuned models.

\subsection{Contributions}

All contributions flow from a single observable---the truncation signature:
\begin{enumerate}
    \item \textbf{Free correctness signal} (\S5): Token count AUC $0.865$ vs.\ $0.561$ for log-probability, requiring no model internals.
    \item \textbf{Two-Type Failure Framework} (\S6): Type~A (truncation: detectable, recoverable) vs.\ Type~B (reasoning: undetectable, irrecoverable) classification.
    \item \textbf{Three-zone budget heuristic} (\S7): A ceiling-rate-based decision heuristic validated across datasets and model families ($r = -0.987$).
    \item \textbf{Deterministic asymmetry} (\S4): The Ratchet Effect---$0/521$ natural-stop recovery under greedy decoding, validated at temperature $= 0.7$---provides the theoretical foundation.
\end{enumerate}

\section{Related Work}

\textbf{Test-Time Compute Scaling.}
Chain-of-thought prompting~\cite{wei2022chain} and self-consistency~\cite{wang2022self} established that additional inference computation can improve reasoning. Snell et al.~\cite{snell2024scaling} proposed optimal allocation of test-time compute, while Sardana et al.~\cite{sardana2024chinchilla} extended scaling laws to inference costs. Recent work on inference scaling~\cite{brown2024large,wu2024inference} examines compute-optimal strategies. These works assume continuous, monotonic scaling. Our findings reveal qualitative phase transitions and a sharp asymmetry in how wrong answers respond to additional compute.

\textbf{Overthinking in Reasoning Models.}
Chen et al.~\cite{chen2024overthinking} identified overthinking in o1-like models: accuracy degrades beyond optimal reasoning length. Subsequent surveys~\cite{stoptoverthinking2025survey} classify approaches as model-based, output-based, or prompt-guided. Budget forcing~\cite{muennighoff2025s1} constrains thinking length. Recent work shows overthinking extends to test-time compute scaling~\cite{morethinkinghurts2026} and proposes statistical early stopping~\cite{statstop2026}. A critical open question is whether overthinking is universal or model-scale-dependent. Our DeepSeek-R1 experiments on one distilled 7B model show no accuracy degradation with increasing budget.

\textbf{Self-Correction and Recovery.}
Huang et al.~\cite{huang2023selfcorrect} showed that LLMs cannot self-correct reasoning errors without external feedback. Subsequent work~\cite{tyen2024llms,madaan2023self} examined conditions under which self-correction helps. Recent evidence suggests that models often determine answers early in generation and produce post-decision explanation~\cite{decideearly2026}, and that reflection after an initial wrong answer provides limited benefit~\cite{firsttry2025}. Our Recovery Rule is a stronger finding: not only can models not \emph{actively} self-correct, they cannot \emph{passively} self-correct either. Even when given more tokens unconditionally, natural-stop wrong answers remain wrong ($0/521$).

\textbf{Efficiency and Verification.}
Process reward models (PRMs)~\cite{lightman2023process} provide step-wise verification but require training data with step-level annotations, as do generative PRMs~\cite{genprm2024,zhang2025generative} and outcome reward models~\cite{cobbe2021gsm8k}. These approaches require an external model trained on task-specific data, whereas our ceiling-rate signal requires only output length. Prior work analyzed the relationship between output length and correctness~\cite{concisethoughts2024} but without distinguishing truncation from reasoning failures---the distinction that makes the signal practical. Concurrent work on inference-time scaling~\cite{qiu2024treebon,setlur2025scaling} examines compute-optimal strategies for repeated sampling and test-time search; our approach is complementary, focusing on per-sample budget decisions based on a free signal. Reasoning methods such as tree-of-thought~\cite{yao2024tree} and graph-of-thought~\cite{besta2024graph} extend chain-of-thought~\cite{wei2022chain} to structured search; our work examines how budget constraints affect the reliability of model outputs at inference time.

\section{Experimental Setup}

\textbf{Models.}
We evaluate 8 models from 3 families: Qwen2.5-\{0.5B, 3B, 7B\}-Instruct (standard instruction-tuned), Gemma-2-\{2B, 9B\}-IT, LLaMA-3-8B-Instruct, and DeepSeek-R1-Distill-Qwen-7B (reasoning model, 7.62B params). We exclude Qwen2.5-1.5B-Instruct as its intermediate scale between 0.5B and 3B provides limited additional discriminative power for our analysis. The three families span different architectures and training paradigms, enabling cross-family validation.

\textbf{Datasets.}
MATH~\cite{hendrycks2021math} (competition mathematics, 200 problems), GSM8K~\cite{cobbe2021gsm8k} (grade-school arithmetic), and HumanEval~\cite{humaneval2021} (Python code generation, 164 problems). All models evaluated on identical subsets for paired comparison.

\textbf{Budgets.}
Standard models: 128--1024 tokens (6 points). Reasoning model: 256--8192 tokens (8 points; virtual budgets extracted from single @8192 generation via token-level truncation; greedy decoding ensures equivalence). Primary experiments use greedy decoding (temperature=0); a sampling validation at temperature $= 0.7$ is reported in \S4.

\textbf{Virtual Budget Extraction.}
For the reasoning model, we generate once at @8192 and truncate at the token level to obtain results at lower budgets. Because we use greedy decoding (temperature=0, top\_k=1), the truncated outputs are identical to what would be produced with the lower budget directly---this is a consequence of autoregressive generation with deterministic decoding. We verified this equivalence on a 10-sample pilot, confirming exact match in all cases. This method cannot be used for the sampling experiments (\S4), which require independent stochastic generation at each budget.

\textbf{Answer Extraction.}
Tiered regex: \texttt{\textbackslash boxed\{\}} notation (highest priority), explicit answer patterns (``the answer is'', ``therefore''), last number in final line. Numerical tolerance $|a-b| < 10^{-6}$ with string matching fallback.

\textbf{Statistical Analysis.}
We use independent sample counting for the Ratchet Effect: each sample is counted once at the most informative budget transition (the transition with the most natural-stop errors for that model). Fisher's exact test for comparing recovery rates between natural-stop and ceiling-hit errors. Clopper-Pearson 95\% CIs for recovery rates. Per-sample bootstrap 95\% CIs for accuracy (10,000 resamples). Cohen's $d$ for effect sizes. All results on 200 samples.

\section{The Ratchet Effect}

We begin with the Ratchet Effect because it provides the theoretical foundation for everything that follows: the token signal (\S5), the two-type framework (\S6), and the budget allocation rule (\S7) all depend on the asymmetry between natural-stop and ceiling-hit error recovery.

\textbf{Definition.}
For a model $M$ at budget $B_1$, define two error types among wrong answers: \emph{ceiling-hit errors} (output length $\geq B_1 - 5$, where the 5-token buffer accounts for end-of-sequence tokens and partial tokenization) and \emph{natural-stop errors} (output length $< B_1 - 5$). At a higher budget $B_2 > B_1$, the \emph{Ratchet Effect} is the observation that ceiling-hit errors recover at a substantial rate while natural-stop errors \emph{never} recover:
$$P(\text{correct at } B_2 \mid \text{wrong \& ceiling-hit at } B_1) \gg 0$$
$$P(\text{correct at } B_2 \mid \text{wrong \& natural-stop at } B_1) = 0$$

\textbf{Empirical Evidence.}
Table~\ref{tab:ratchet} presents recovery analysis using \emph{independent sample counting}: each sample is counted once at its most informative transition (the budget pair with the most natural-stop errors for that model). The point estimate (0\% natural-stop recovery) is identical at every transition; the choice of transition affects only the confidence interval width.

\begin{table}[t]
\centering
\caption{The Ratchet Effect across 3 model families and 3 datasets. \textbf{NatW}: unique natural-stop wrong; \textbf{NatR}: recovered; \textbf{CeilW}: ceiling-hit wrong; \textbf{CeilR}: recovered. Combined: $0/521$ natural-stop vs.\ $438/804$ ($54.5\%$) ceiling-hit recovery.}
\label{tab:ratchet}
\begin{tabular}{lllcccccc}
\toprule
\textbf{Family} & \textbf{Model} & \textbf{Dataset} & \textbf{NatW} & \textbf{NatR} & \textbf{CeilW} & \textbf{CeilR} & \textbf{CeilR\%} \\
\midrule
\multirow{5}{*}{Qwen2.5} & 0.5B & MATH & 22 & 0 & 127 & 37 & 29.1\% \\
 & 3B & MATH & 14 & 0 & 115 & 72 & 62.6\% \\
 & 7B & MATH & 20 & 0 & 101 & 66 & 65.3\% \\
 & 3B & GSM8K & 9 & 0 & 108 & 84 & 77.8\% \\
 & 7B & GSM8K & 12 & 0 & 93 & 74 & 79.6\% \\
\midrule
\multirow{4}{*}{Gemma-2} & 2B-IT & MATH & 70 & 0 & 21 & 0 & 0.0\% \\
 & 2B-IT & GSM8K & 69 & 0 & 18 & 3 & 16.7\% \\
 & 9B-IT & MATH & 47 & 0 & 18 & 4 & 22.2\% \\
 & 9B-IT & GSM8K & 26 & 0 & 13 & 7 & 53.8\% \\
\midrule
LLaMA-3 & 8B & GSM8K & 95 & 0 & 1 & 0 & ---$^{\ddagger}$ \\
\midrule
\multirow{1}{*}{Qwen2.5} & 7B & HumanEval & 99 & 0 & 35 & 1 & 2.9\% \\
\midrule
R1-Distill & 7B & MATH & 38 & 0 & 2 & 2 & 100\% \\
 & 7B & GSM8K & 0 & 0 & 152 & 88 & 57.9\% \\
\midrule
\multicolumn{3}{l}{\textbf{Total (13 conditions)}} & \textbf{521} & \textbf{0} & \textbf{804} & \textbf{438} & \textbf{54.5\%} \\
\bottomrule
\end{tabular}
\vspace{1mm}

\raggedright\small $^{\ddagger}$LLaMA-3-8B produces very short outputs (avg.\ $\sim$3 tokens), yielding only 1 ceiling-hit sample; recovery rate not estimable. R1-Distill-7B on MATH shows 2/2 ceiling recovery; with only 2 ceiling-hit samples, this rate is not reliably estimable. R1-Distill-7B on GSM8K has 0 natural-stop errors at @256 (99\% truncation rate), providing no Type~B samples at this transition. Qwen2.5-7B on HumanEval uses the 768$\to$1024 transition (most natural-stop errors); ceiling-hit recovery is 2.9\%, reflecting the model's limited code generation capability at this scale. Gemma-2-2B ceiling recovery is 0\%, reflecting limited reasoning capability at this scale.
\end{table}

\begin{figure}[t]
\centering
\includegraphics[width=0.85\columnwidth]{figures/fig_ratchet_effect.pdf}
\caption{The Ratchet Effect: natural-stop errors (gray) never recover while ceiling-hit errors (colored) recover at 0--80\% across 3 model families. Each bar represents one model-dataset condition.}
\label{fig:ratchet}
\end{figure}

\textbf{Statistical Significance.}
We apply Fisher's exact test to the $2 \times 2$ contingency table (natural-stop: 0 recovered of 521; ceiling-hit: 438 recovered of 804), yielding $p < 10^{-99}$ (odds ratio $= \infty$). This tests the genuine empirical question---whether ceiling-hit recovery strictly exceeds natural-stop recovery---rather than the deterministic natural-stop rate itself. The \textbf{Clopper-Pearson 95\% CI} for the natural-stop recovery rate is $[0, 0.9\%]$, though we note this is a confidence interval around a mathematical certainty under deterministic decoding.

\textbf{Robustness to Analysis Window and Cross-Family Validation.}
Across all 15 available model-dataset-transition combinations (128$\to$256, 256$\to$512, 512$\to$1024), the natural-stop recovery rate remains exactly 0\% in every case. DeepSeek-R1-Distill-7B shows the same pattern: $0/38$ natural-stop recoveries across all budget transitions. The Ratchet Effect is not driven by any single model, family, or dataset: it holds across Qwen2.5, Gemma-2, and LLaMA-3 architectures, for small models (0.5B) and large models (9B), standard and reasoning models, on math (MATH, GSM8K) and code (HumanEval) tasks. LLaMA-3-8B produces very short outputs ($\sim$3 tokens), contributing 95 of 521 natural-stop errors. Excluding it entirely, the result remains $0/426$ (Fisher's exact $p < 10^{-86}$).

\textbf{Interpretation.}
Under deterministic decoding, natural-stop non-recovery is a mathematical certainty of autoregressive generation, not an empirical estimate: the model has completed its reasoning and produced a final answer, so extending the budget reproduces the identical output. The contribution is not this qualitative observation (which is expected) but its practical consequence: (1) this certainty enables a zero-overhead heuristic---never extend the budget for natural-stop outputs; (2) ceiling-hit recovery varies from $0\%$ (Gemma-2B) to $80\%$ (Qwen2.5-7B), revealing a model-capability gradient that is genuinely empirical. The practical value lies in the asymmetric gap between deterministic Type~B non-recovery and empirical Type~A recovery.

\textbf{Sampling Validation (Temperature $= 0.7$).}
To test whether the practical asymmetry extends beyond greedy decoding, we ran Qwen2.5-7B-Instruct on MATH 200 samples with temperature $= 0.7$, top\_p $= 0.9$ across 3 seeds at budgets 256 and 512. Under same-seed sampling (deterministic), natural-stop recovery remains exactly $0/N$ across all 3 seeds (seed42: $0/18$; seed123: $0/19$; seed456: $0/17$), confirming the Ratchet Effect extends to any deterministic decoding. Under cross-seed sampling (independent stochastic generations), the budget asymmetry persists: natural-stop errors show 26--44\% fresh-attempt correctness while ceiling-hit errors reach 54--61\%. A control experiment at the \emph{same} budget (@256$\to$@256 with different seeds) disentangles two effects: sampling variation helps both error types equally, but \emph{budget increase selectively helps ceiling-hit errors} (+44pp vs.\ +2--13pp for natural-stop, see Appendix). The Type~A/Type~B distinction is therefore not an artifact of greedy decoding.

\section{The Truncation Signal}

\textbf{Motivation.}
If wrong answers fall into two types with radically different properties, can we \emph{detect} which type a given answer is? The answer for Type~A is trivially yes: checking whether output length equals the budget directly identifies truncation. The practical question is whether this trivial detector suffices at typical budgets, and how it compares to the standard approach of using model confidence (log-probability).

\textbf{Setup.}
For each model and budget, we compute AUC-ROC for distinguishing correct from wrong answers using: (a) output token count, and (b) mean log-probability of the output.

\begin{table}[t]
\centering
\caption{Token count vs.\ log-probability for detecting wrong answers at 256 tokens. Token count dominates on all 8 conditions. Signal independence: all $|r| < 0.2$. Cohen's $d$ measures the standardized effect size between correct and wrong answer distributions for each signal.}
\label{tab:signal}
\begin{tabular}{llcccc}
\toprule
\textbf{Model} & \textbf{Dataset} & \textbf{Token AUC} & \textbf{Logprob AUC} & \textbf{Ratio} & \textbf{Cohen's $d$} \\
\midrule
Qwen2.5-7B & MATH & \textbf{0.865} & 0.561 & 1.54$\times$ & 1.43 \\
Qwen2.5-3B & MATH & \textbf{0.840} & 0.543 & 1.55$\times$ & 1.35 \\
Qwen2.5-0.5B & MATH & \textbf{0.794} & 0.736 & 1.08$\times$ & 0.42 \\
Qwen2.5-7B & GSM8K & \textbf{0.880} & 0.465 & 1.89$\times$ & 1.43 \\
Gemma-2-9B & MATH & \textbf{0.731} & 0.584 & 1.25$\times$ & 0.61 \\
Gemma-2-9B & GSM8K & \textbf{0.699} & 0.538 & 1.30$\times$ & 0.52 \\
LLaMA-3-8B & GSM8K & \textbf{0.713} & 0.601 & 1.19$\times$ & 0.84 \\
Qwen2.5-7B & HumanEval & \textbf{0.884} & 0.650 & 1.36$\times$ & 0.92 \\
\bottomrule
\end{tabular}
\end{table}

\begin{figure}[t]
\centering
\includegraphics[width=\columnwidth]{figures/fig_token_signal_roc.pdf}
\caption{ROC curves for token count vs.\ log-probability as wrong-answer detectors on Qwen2.5-7B. Token count (red) substantially outperforms log-probability (blue) on both MATH (AUC 0.865 vs.\ 0.561) and GSM8K (AUC 0.880 vs.\ 0.465).}
\label{fig:signal_roc}
\end{figure}

\textbf{Results.}
Table~\ref{tab:signal} and Figure~\ref{fig:signal_roc} show that token count substantially outperforms log-probability on deployed models (7B+). On Qwen2.5-7B MATH, token AUC ($0.865$) exceeds log-prob AUC ($0.561$) by $1.54\times$. On GSM8K, the gap widens to $1.89\times$ ($0.880$ vs.\ $0.465$). The signal transfers cross-family: Gemma-2-9B achieves token AUC of $0.699$--$0.731$, lower than Qwen2.5-7B but well above random, consistent with Gemma's lower ceiling rates at 256 tokens. Log-probability provides essentially no signal for detecting wrong answers at this scale. Across all 8 conditions, token count dominates by $1.08$--$1.89\times$. The signal extends beyond mathematical reasoning: on HumanEval (code generation), token AUC reaches $0.884$ at 256 tokens---even stronger than on math, reflecting the structured nature of code outputs where truncation more reliably indicates incompleteness.

\textbf{Why Token Count Works.}
The token signal works because truncation failures dominate the error set at typical budgets. At 256 tokens, the 7B model is truncated in 55\% of responses; among wrong answers, the vast majority are ceiling-hit errors. Since ceiling-hit errors have token count $\approx$ budget (trivially distinguishable from correct answers), the signal is strong. Critically, the two signals are independent ($|r| < 0.2$ across all conditions), confirming they capture different information. We emphasize that this is not a ``discovered'' signal but a quantification of practical value: the trivial check of output length suffices because Type~A errors dominate, and Type~B errors are irrecoverable regardless.

\textbf{The Signal Fails for Reasoning Errors.}
Type~B errors (natural-stop wrong) are indistinguishable from correct answers: the natural-wrong vs.\ correct AUC is $\approx 0.5$---random chance. The token signal is blind to reasoning quality. This limitation is inherent: output length cannot assess reasoning correctness. However, this limitation is practically benign because reasoning errors \emph{never recover anyway} (Ratchet Effect), so detecting them provides no actionable information for budget allocation.

\section{Two-Type Failure Framework}

The Ratchet Effect and Truncation Signal converge on a single framework (Table~\ref{tab:theory}):

\begin{table}[h]
\centering
\caption{Two-Type Failure Framework: Two failure types with distinct properties.}
\label{tab:theory}
\begin{tabular}{lcc}
\toprule
& \textbf{Type~A: Truncation} & \textbf{Type~B: Reasoning} \\
\midrule
Cause & Budget too short & Model reasoning failed \\
Token count & $\approx$ budget & $<$ budget \\
Detectable (AUC) & $\approx 1.0$ & $\approx 0.5$ (random) \\
Recoverable & 0--80\% & 0\% ($0/521$) \\
Token signal & \textbf{Strong} & \textbf{Blind} \\
\bottomrule
\end{tabular}
\end{table}

\textbf{Proof of Separability.}
The two failure types occupy disjoint regions of the token-count feature space. On 7B at 256 tokens: ceiling-wrong vs.\ correct AUC $= 0.865$; natural-wrong vs.\ correct AUC $= 0.168$ (below 0.5 because natural-wrong answers stop at short lengths while correct answers often run longer, completing their reasoning). On 3B at 256 tokens: ceiling-wrong AUC $= 0.815$; natural-wrong AUC $= 0.168$. On 0.5B at 256 tokens: ceiling-wrong AUC $= 0.794$. At 512 tokens for 7B and 3B: ceiling-wrong AUC $> 0.99$ (near-perfect separation); natural-wrong AUC $\approx 0.5$ (random). The separability holds consistently across model scales.

\textbf{Key Insight.}
The token signal works not because it detects reasoning quality, but because truncation failures \emph{dominate} the wrong set at practical budgets. And that dominance is sufficient, because reasoning failures never recover anyway---there is no practical need to detect them.

\section{The Token Budget Spectrum}

\textbf{Ceiling Rate as Diagnostic.}
The fraction of outputs hitting the budget ceiling is a universal indicator of model-budget alignment. Across 7B models (standard + reasoning, 6 budgets each), ceiling rate correlates with accuracy at Pearson $r = -0.987$; within each model family, the relationship is monotonic (per-family Spearman $\rho = -1.000$). This motivates a three-zone heuristic:

\begin{itemize}
    \item \textbf{TRUNCATION zone} (ceiling rate $> 50\%$): Budget severely constrains the model. Extending the budget has high returns.
    \item \textbf{INFLECTION zone} (ceiling rate $10$--$50\%$): Budget partially constrains the model. Cautious extension recommended.
    \item \textbf{EFFICIENT zone} (ceiling rate $< 10\%$): Budget is sufficient. Further extension has diminishing returns.
\end{itemize}

These thresholds are empirically grounded: across our data, ceiling rates above 50\% always correspond to the TRUNCATION zone (accuracy gain $>$ 10pp per budget doubling), while below 10\% the gain is $<$ 1pp. The thresholds are robust to moderate variation: shifting to 40\%/60\% or 5\%/15\% produces the same zone assignments for all conditions in Table~\ref{tab:spectrum}.

\textbf{Standard vs.\ Reasoning Models.}
Table~\ref{tab:spectrum} and Figure~\ref{fig:spectrum} compare the budget spectrum of standard Qwen2.5-7B-Instruct and DeepSeek-R1-Distill-7B.

\begin{table}[t]
\centering
\caption{Token Budget Spectrum: Standard vs.\ Reasoning model on MATH / GSM8K. Both reach similar peak accuracy ($80.5\%$ / $81.0\%$), but R1 requires $2.0\times$ more tokens ($@1024$ vs.\ $@512$) to reach the EFFICIENT zone. R1 values show MATH / GSM8K; R1 at 384 and 768 computed via virtual budget extraction from @8192. Zone assignments are consistent across datasets.}
\label{tab:spectrum}
\begin{tabular}{lcccccc}
\toprule
& \multicolumn{3}{c}{\textbf{Qwen2.5-7B (Standard)}} & \multicolumn{3}{c}{\textbf{R1-Distill-7B (Reasoning)}} \\
\cmidrule(lr){2-4} \cmidrule(lr){5-7}
\textbf{Budget} & Acc & Ceiling & Zone & Acc & Ceiling & Zone \\
\midrule
128 & 5.0\% & 99.5\% & TRUNC & ---$^{\S}$ & ---$^{\S}$ & ---$^{\S}$ \\
256 & 53.5\% & 55.5\% & TRUNC & 12.5/24.0\% & 96/99\% & TRUNC \\
384 & 76.0\% & 11.0\% & INFLECT & 35.0/36.5\% & 69/67\% & TRUNC \\
512 & 78.5\% & 2.5\% & EFF & 63.5/67.0\% & 36/29\% & INFLECT \\
768 & 80.5\% & 0.0\% & EFF & 78.0/79.0\% & 8/5\% & EFF \\
1024 & 80.5\% & 0.0\% & EFF & 80.0/80.5\% & 3/2\% & EFF \\
2048 & --- & --- & --- & 80.5/81.0\% & 0.5/0.5\% & EFF \\
8192 & --- & --- & --- & 80.5/81.0\% & 0.5/0\% & EFF \\
\bottomrule
\end{tabular}
\vspace{1mm}

\raggedright\small $^{\S}$R1-Distill-7B not evaluated at @128 (reasoning models require sufficient budget for the thinking phase; 128 tokens produce near-universal truncation with no meaningful answer).
\end{table}

\begin{figure}[t]
\centering
\includegraphics[width=0.9\columnwidth]{figures/fig_budget_spectrum_v3.pdf}
\caption{Token Budget Spectrum for standard (blue) and reasoning (red) 7B models on MATH. Both converge to $>$80\% peak accuracy (80.5\% on MATH, 81.0\% on GSM8K for R1). The reasoning model is right-shifted: it enters the EFFICIENT zone at @1024 vs.\ @512 for the standard model.}
\label{fig:spectrum}
\end{figure}

Three findings emerge:

\textbf{1. Peak convergence.} Both models reach similar peak accuracy (80.5\% on MATH; 81.0\% for R1 on GSM8K)---the performance ceiling is determined by task difficulty, not model architecture. We note that on 200 samples, a 0.5\% difference corresponds to one problem, so this convergence should be interpreted within sampling noise.

\textbf{2. Reasoning model right-shift.} R1 needs @1024 tokens to enter the EFFICIENT zone, while the standard model enters at @512. The reasoning model's internal thinking tokens consume budget before answer generation, shifting its truncation boundary rightward by $2.0\times$.

\textbf{3. Monotonic accuracy increase.} R1 accuracy is monotonically non-decreasing across all 6 evaluated budgets from @256 (12.5/24.0\%) through @384 (35.0/36.5\%), @512 (63.5/67.0\%), @768 (78.0/79.0\%), @1024 (80.0/80.5\%) to @8192 (80.5/81.0\%) on MATH/GSM8K. At @8192, the model uses an average of only 536 tokens. No accuracy degradation is observed at any budget, though we caution that this is a single distilled model and may not generalize to larger, RL-trained reasoning models where overthinking has been documented~\cite{chen2024overthinking}.

\textbf{Implications for Overthinking Research.}
Recent work~\cite{chen2024overthinking,stoptoverthinking2025survey} identifies overthinking as a critical challenge for reasoning models. Our findings from one distilled 7B model suggest that, at this scale, under-allocation warrants attention alongside overthinking. Whether this extends to larger reasoning models remains an open question.

\section{Application: Prospective Budget Allocation}

The Ratchet Effect suggests a simple adaptive strategy: run at initial budget $B_1$ (e.g., 256 tokens); if the output hits the ceiling, extend to $B_2$; if the output stops naturally, do not extend. This heuristic uses only the observable ceiling-hit status---no ground truth or trained verifier required. The logic is: ceiling-hit outputs have a $54.5\%$ recovery rate (extending often helps), while natural-stop outputs have $0/521$ recovery (extending never helps). Note that this extends even correct ceiling-hit answers; the cost of occasionally re-generating a correct answer is outweighed by the $54.5\%$ chance of recovering a wrong one. The key advantage is \emph{zero overhead}: no model internals, no difficulty estimator---only the trivially observable output length. The three-zone heuristic (TRUNCATION: extend; EFFICIENT: stop; INFLECTION: cautious extension) refines this binary decision based on ceiling rate magnitude.

\textbf{Prospective Validation.}
To test whether the heuristic produces correct decisions on unseen data, we run Qwen2.5-7B on 50 held-out MATH problems (Level 1--3, not used in any analysis above): Phase~1 generates all 50 at @256; Phase~2 extends only ceiling-hit outputs to @512. The model's ceiling rate at @256 is 6.0\% (3/50), placing it firmly in the EFFICIENT zone---the heuristic correctly recommends no extension. Adaptive accuracy matches the full-budget baseline (92.0\% vs.\ 92.0\%, $n=50$), while total compute drops by 86.4\%. Natural-stop recovery is 0/4, consistent with the Ratchet Effect. We note that this held-out set consists of easier problems (Level 1--3) where the model is already competent; a prospective evaluation on harder problems (Level 4--5) where the model enters the TRUNCATION zone would provide a stronger test of the compute--accuracy tradeoff, and we leave this to future work with access to harder held-out splits.

\section{Discussion}

\textbf{Generalizability.}
The Ratchet Effect ($0/521$) holds across three model families (Qwen2.5, Gemma-2, LLaMA-3) and three tasks: mathematical reasoning (MATH), arithmetic (GSM8K), and code generation (HumanEval). The ceiling-hit recovery rate varies substantially---Qwen2.5-7B achieves 65--80\% on math but only 2.9\% on HumanEval, reflecting differences in task difficulty and model capability. Gemma-2-2B represents a boundary case: its ceiling-hit errors are correctly classified as Type~A (detectable) yet show low recovery, revealing that Type~A recoverability is \emph{model-dependent}---truncation removes the budget constraint, but the model must still produce the correct answer.

\textbf{When Recovery Is Low.}
The framework provides the most actionable budget-saving information when models are competent enough to benefit from extension (high ceiling-hit recovery). When models lack task capability---as with Qwen2.5-7B on HumanEval (2.9\% recovery) or Gemma-2-2B on MATH (0\% recovery)---the heuristic correctly identifies truncation, but extension simply produces a different wrong answer. This is not a failure of the detection mechanism but a correct diagnosis: the bottleneck is model capability, not budget allocation. In practice, low ceiling-hit recovery combined with high truncation rate signals that the model-dataset pair is mismatched at this scale, and budget extension alone cannot compensate.

\textbf{Statistical Power.}
The universal claim rests on cross-condition consistency (13/13 conditions show 0\% natural-stop recovery) rather than per-condition sample size. Pooling gives $n=521$ for natural-stop and $n=804$ for ceiling-hit errors (Fisher's exact $p < 10^{-99}$). All effect sizes (0\% vs.\ 54.5\% recovery; AUC 0.865 vs.\ 0.561; $r = -0.987$) yield power $= 1.000$. Excluding LLaMA-3-8B: $0/426$ (Fisher's exact $p < 10^{-86}$).

\textbf{Connection to Self-Correction.}
Huang et al.~\cite{huang2023selfcorrect} showed that LLMs cannot actively self-correct without external feedback. Our finding is strictly stronger: models cannot \emph{passively} self-correct either---even given more tokens unconditionally, natural-stop errors remain wrong ($0/521$).

\textbf{Limitations.}
Primary experiments use greedy decoding; the sampling validation (\S4) at temperature $= 0.7$ confirms the asymmetry extends to stochastic generation, but a systematic study across decoding strategies (nucleus sampling, diverse beam search) and temperature settings is needed. Tasks cover mathematical reasoning (MATH, GSM8K) and code generation (HumanEval); validation on commonsense reasoning and open-ended tasks would further strengthen generalizability. The prospective validation (\S8) uses 50 held-out MATH problems at Level 1--3 difficulty; a stronger test would evaluate harder problems where the model enters the TRUNCATION zone. The R1-Distill-7B result is a single model on two datasets (MATH, GSM8K). LLaMA-3-8B produces very short outputs ($\sim$3 tokens), reflecting instruction-following limitations; the main result holds without it. The three-zone thresholds (10\%, 50\%) are empirically grounded but should be validated on broader task distributions.

\section{Conclusion}

We show that output token count---a trivially observable, free signal---suffices for test-time compute budget optimization. The signal works because model failures split into two types: truncation failures (detectable via token count $\approx$ budget, recoverable 0--80\% with more tokens) and reasoning failures (undetectable, irrecoverable under deterministic decoding). This Two-Type Failure Framework yields a three-zone budget allocation heuristic based on ceiling rate (TRUNCATION / INFLECTION / EFFICIENT) that requires no verifiers, training, or ground truth. The deterministic non-recovery of reasoning failures ($0/521$ across 3 families and 3 tasks) provides the theoretical foundation; the practical value is the token signal itself (AUC $0.865$ vs.\ $0.561$ for log-probability). The rule extends to reasoning models: DeepSeek-R1-Distill-7B shows a $2.0\times$ right-shift in the truncation boundary but reaches the same accuracy ceiling (80.5--81.0\%), providing initial evidence that the framework applies beyond standard instruction-tuned models.

\section*{Ethics Statement}

Increased inference computation raises energy concerns; our adaptive algorithm reduces unnecessary compute. We release our evaluation framework for reproducible research.

\bibliographystyle{neurips_2024}
\begin{thebibliography}{30}

\bibitem{wei2022chain}
Jason Wei, Xuezhi Wang, Dale Schuurmans, Maarten Bosma, Fei Xia, Ed Chi, Quoc V Le, Denny Zhou, et al.
\newblock Chain-of-thought prompting elicits reasoning in large language models.
\newblock \emph{NeurIPS}, 2022.

\bibitem{wang2022self}
Xuezhi Wang, Jason Wei, Dale Schuurmans, Quoc Le, Ed Chi, Sharan Narang, Aitor Chowdhery, and Denny Zhou.
\newblock Self-consistency improves chain of thought reasoning in language models.
\newblock \emph{ICLR}, 2023.

\bibitem{snell2024scaling}
Charlie Snell, Jaehoon Lee, Kelvin Xu, and Aviral Kumar.
\newblock Scaling LLM test-time compute optimally can be more effective than scaling model parameters.
\newblock \emph{arXiv preprint arXiv:2408.03314}, 2024.

\bibitem{sardana2024chinchilla}
Nishkam Ravi Sardana, Jacob Portes, Soumya Doubov, and Jonathan Frankle.
\newblock Beyond chinchilla-optimal: Accounting for inference in language model scaling laws.
\newblock \emph{arXiv preprint arXiv:2401.00448}, 2024.

\bibitem{brown2024large}
Bradley Brown, Jordan Juravsky, Ryan Ehrlich, et al.
\newblock Large language monkeys: Scaling inference compute with repeated sampling.
\newblock \emph{arXiv preprint arXiv:2407.21787}, 2024.

\bibitem{wu2024inference}
Yangzhen Wu, Zhiqing Sun, Shanda Li, Sean Welleck, and Yiming Yang.
\newblock Inference scaling laws: An empirical analysis of compute-optimal inference for problem-solving with language models.
\newblock \emph{arXiv preprint arXiv:2408.00724}, 2024.

\bibitem{chen2024overthinking}
Xingyu Chen, Jiahao Xu, Tian Liang, Zhiwei He, Ruisheng Peng, and Yu Qiao.
\newblock Do NOT think that much for 2+3=? On the overthinking of o1-like LLMs.
\newblock \emph{arXiv preprint arXiv:2412.21187}, 2024.

\bibitem{stoptoverthinking2025survey}
Yang Sui, Yu-Neng Chuang, Guanchu Wang, Jiamu Zhang, Tianyi Zhang, et al.
\newblock Stop overthinking: A survey on efficient reasoning for large language models.
\newblock \emph{TMLR}, 2025.

\bibitem{huang2023selfcorrect}
Jie Huang, Xinyun Chen, Swaroop Mishra, Huaixiu Steven Zheng, Adams Wei Yu, Xinying Song, and Denny Zhou.
\newblock Large language models cannot self-correct reasoning yet.
\newblock \emph{ICLR}, 2024.

\bibitem{tyen2024llms}
Gladys Tyen, Hassan Mansoor, Peter Chen, Tony Mak, and Victor Carbune.
\newblock LLMs cannot find reasoning errors, but can correct them.
\newblock \emph{ICLR}, 2024.

\bibitem{madaan2023self}
Aman Madaan, Niket Tandon, Prakhar Gupta, et al.
\newblock Self-refine: Iterative refinement with self-feedback.
\newblock \emph{NeurIPS}, 2023.

\bibitem{muennighoff2025s1}
Niklas Muennighoff, Zitong Yang, Weijia Shi, et al.
\newblock s1: Simple test-time scaling.
\newblock \emph{arXiv preprint arXiv:2501.19393}, 2025.

\bibitem{lightman2023process}
Hunter Lightman, Vineet Kosaraju, Yunda Li, et al.
\newblock Let's verify step by step.
\newblock \emph{NeurIPS}, 2023.

\bibitem{genprm2024}
Qihao Zhao, Yang Liu, Qichao Zhang, et al.
\newblock GenPRM: Scaling test-time compute of process reward models via generative reasoning.
\newblock \emph{arXiv preprint arXiv:2504.00891}, 2025.

\bibitem{zhang2025generative}
Lunjun Zhang, Arian Hosseini, Hritik Bansal, Mehran Kazemi, Aviral Kumar, and Rishabh Agarwal.
\newblock Generative verifiers: Reward modeling as next-token prediction.
\newblock \emph{arXiv preprint arXiv:2408.15240}, 2024.

\bibitem{cobbe2021gsm8k}
Karl Cobbe, Vineet Kosaraju, Mohammad Bavarian, et al.
\newblock Training verifiers to solve math word problems.
\newblock \emph{arXiv preprint arXiv:2110.14168}, 2021.

\bibitem{qiu2024treebon}
Jian Qiu, Yifu Lu, Yizhe Zeng, Jie Guo, Jian Geng, Hao Wang, Kai Huang, Yanhua Wu, and Meng Wang.
\newblock TreeBoN: Enhancing inference-time alignment with speculative tree-search best-of-N sampling.
\newblock \emph{arXiv preprint arXiv:2410.16033}, 2024.

\bibitem{setlur2025scaling}
Amrith Setlur, Nived Rajaraman, Sergey Levine, and Aviral Kumar.
\newblock Scaling test-time compute without verification or RL is suboptimal.
\newblock \emph{arXiv preprint arXiv:2502.12118}, 2025.

\bibitem{humaneval2021}
Mark Chen, Jerry Tworek, Heewoo Jun, et al.
\newblock Evaluating large language models trained on code.
\newblock \emph{arXiv preprint arXiv:2107.03374}, 2021.

\bibitem{hendrycks2021math}
Dan Hendrycks, Collin Burns, Saurav Kadavath, et al.
\newblock Measuring mathematical problem solving with the MATH dataset.
\newblock \emph{NeurIPS}, 2021.

\bibitem{yao2024tree}
Shunyu Yao, Dian Yu, Jeffrey Zhao, et al.
\newblock Tree of thoughts: Deliberate problem solving with large language models.
\newblock \emph{NeurIPS}, 2023.

\bibitem{besta2024graph}
Maciej Besta, Nils Blach, Ales Kubicek, et al.
\newblock Graph of thoughts: Solving elaborate problems with large language models.
\newblock \emph{AAAI}, 2024.

\bibitem{guo2025deepseek}
Daya Guo, Dejian Yang, et al.
\newblock DeepSeek-R1: Incentivizing reasoning capability in LLMs via reinforcement learning.
\newblock \emph{arXiv preprint arXiv:2501.12948}, 2025.

\bibitem{firsttry2025}
Georgios Chochlakis, Kristina Lerman, and Aarash Feizi.
\newblock First try matters: Revisiting the role of reflection in reasoning models.
\newblock \emph{arXiv preprint arXiv:2510.08308}, 2025.

\bibitem{concisethoughts2024}
Maxime Lupion, Romain Bret, and Iacer Calixto.
\newblock Concise thoughts: Impact of output length on LLM reasoning and cost.
\newblock \emph{arXiv preprint arXiv:2407.19825}, 2024.

\bibitem{decideearly2026}
Soroush Moazed, Arian Hosseini, and Daniel Hsu.
\newblock Large language models decide early and explain later.
\newblock \emph{arXiv preprint arXiv:2604.22266}, 2026.

\bibitem{morethinkinghurts2026}
Vijay Viswanathan, Chen Zhao, and Graham Neubig.
\newblock When more thinking hurts: Overthinking in LLM test-time compute scaling.
\newblock \emph{arXiv preprint arXiv:2604.10739}, 2026.

\bibitem{statstop2026}
Jiahao Xu, Tian Liang, Zhiwei He, Ruisheng Peng, and Yu Qiao.
\newblock Statistical early stopping for reasoning models.
\newblock \emph{arXiv preprint arXiv:2602.13935}, 2026.

\end{thebibliography}

\section*{Appendix: Sampling Validation Details}

Table~\ref{tab:sampling} provides the full results of the sampling experiment at temperature $= 0.7$.

\begin{table}[h]
\centering
\caption{Sampling validation at temperature $= 0.7$ (Qwen2.5-7B-Instruct, MATH 200). Same-seed: deterministic recovery (Ratchet). Cross-seed/Control: fresh-attempt correctness, denoted NatOK/CeilOK (different seed's independent generation happens to be correct).}
\label{tab:sampling}
\begin{tabular}{llcccc}
\toprule
\textbf{Test} & \textbf{Seed(s)} & \textbf{NatW} & \textbf{NatOK} & \textbf{CeilW} & \textbf{CeilOK} \\
\midrule
\multirow{3}{*}{Same-seed} & 42/42 & 18 & 0 & 121 & 71 \\
 & 123/123 & 19 & 0 & 123 & 71 \\
 & 456/456 & 17 & 0 & 126 & 68 \\
\midrule
\multirow{6}{*}{Cross-seed} & 42$\to$123 & 18 & 6 & 121 & 69 \\
 & 42$\to$456 & 18 & 8 & 121 & 66 \\
 & 123$\to$42 & 19 & 8 & 123 & 75 \\
 & 123$\to$456 & 19 & 5 & 123 & 69 \\
 & 456$\to$42 & 17 & 7 & 126 & 76 \\
 & 456$\to$123 & 17 & 5 & 126 & 71 \\
\midrule
\multirow{6}{*}{Control @256} & 42$\to$123 & 18 & 6 & 121 & 10 \\
 & 42$\to$456 & 18 & 5 & 121 & 14 \\
 & 123$\to$42 & 19 & 6 & 123 & 13 \\
 & 123$\to$456 & 19 & 5 & 123 & 12 \\
 & 456$\to$42 & 17 & 7 & 126 & 16 \\
 & 456$\to$123 & 17 & 4 & 126 & 14 \\
\bottomrule
\end{tabular}
\end{table}

Cross-seed fresh-attempt correctness for natural-stop errors (26--44\%, pooled 36.1\%) is comparable to the control (24--41\%, pooled 32.4\%), indicating most of this improvement comes from sampling variation rather than budget increase. In contrast, ceiling-hit correctness rises from 8--13\% (control, pooled 10.8\%) to 54--61\% (cross-budget, pooled 57.6\%), a +44pp gain attributable to budget extension.

\end{document}


